%!TEX Root = ../main.tex
\chapter{Chain Conditions}
In this chapter, we are going to talk about
chain conditions, which are conditions
on rings / modules that makes various objects finite.
There are two similar conditions, 
Noetherian, which is about ascending chain,
and Artinian, which is about descending chain.
We first discuss these two conditions for modules
in \Cref{sec:chain_modules}.
Later, we will specialize to Noetherian rings
(\Cref{sec:noetherian})
and Artinian rings (\Cref{sec:artinian}).
It will turns out that while Noetherian
rings is a fundamental finiteness condition that makes
rings reasonable, Artinian rings much more specialize
and looks like ``points'' (i.e., has dimension $0$).

\section{Noetherian and Artinian Modules}
\label{sec:chain_modules}
\subsection{Definitions and Basic Examples}
There are two main types of chain condition
that we will discuss.
\begin{definition}[Noetherian and Artininan modules]
    Let $R$ be a ring, and $M$ be an $R$-module.
    \begin{itemize}
    \item $M$ is a \vocab{Noetherian module} if it satisfies
	the \vocab{ascending chain condition (ACC)}:
	there does not exist a chain of submodules
	$$M_1\subsetneq M_2\subsetneq M_3\subsetneq \dots
	\subset M.$$
	(i.e., any ascending chain stops).
    \item $M$ is a \vocab{Artinian module} if it satisfies
	the \vocab{descending chain condition (DCC)}:
	there does not exist a chain of submodules
	$$M\supset M_1\supsetneq M_2\supsetneq M_3
	\supsetneq\dots$$
    \end{itemize}
    We say that $R$ is a \vocab{Noetherian ring} 
    (resp. \vocab{Artinian ring})
    if $R$ is Noetherian (resp. Artinian) as an $R$-module.
\end{definition}
Now, let's do a basic example.
\begin{example}
  We have the following examples.
\begin{itemize}
    \item Finite abelian groups (as a $\mathbb Z$-module)
	are both Noetherian and Artinian
	because they are finite.
    \item $R=\mathbb Z/(n)$ is both
	Noetherian and Artinian for the same reason.
    \item PID's are Noetherian rings:
	if there is a chain
	$(a_1)\subsetneq (a_2)\subsetneq\dots,$
	then $(a_1)\cup (a_2)\cup \dots$ is an ideal 
        generated by element $(a)$.
        Then, $a\in (a_i)$ for some $i$, forcing 
        $(a_i) = (a_{i+1}) = \dots = (a)$, so the chain
        is eventually constant.

    Consequently, $\ZZ$ and $\CC[x]$ are Noetherian.
    \item $\mathbb Z$ is not Artinian ring:
	$\mathbb Z\supset (2)\supset (4)\supset (8)\supset\dots$.
    \item $R=\mathbb Z$ and $M=\mathbb Z[1/p]
	/\mathbb Z$ is Artinian but not Noetherian.
	Any submodule of $M$ is in form
	$M_n = \frac{1}{p^n}\mathbb Z/\mathbb Z$,
	so $0=M_0\subset M_1\subset M_2\subset\dots$.
	Therefore, this module is not Noetherian.
	It is Artinian because any descending chain
	must start somewhere, hence finite.
    \item $\mathbb C[x_1,x_2,\dots]$
	is neither Noetherian nor Artinian ring.
\end{itemize}
\end{example}
We will see later that Noetherian rings
is a more fundamental concept of commutative algebra:
almost all examples you would care about is Noetherian.
On the other hand, Artinian ring is an extremely
special kind of rings that we will describe later.
%
%
%
%
%
\subsection{Properties of Noetherian and Artinian Modules}
\begin{proposition}
    \label{prop:noetherian_module_fin_gen}
    An $R$-module $M$ is Noetherian if and only if
    every submodule is finitely-generated.
\end{proposition}
\begin{proof}
    ($\Rightarrow$) If $N\subset M$ is not finitely generated,
    pick $x_1\in N$, $x_2\notin\Span(x_1)$,
    $x_3\notin\Span(x_1,x_2)$, \dots.
    Then, one has
    $$Rx_1\subsetneq Rx_1+Rx_2\subsetneq
     Rx_1+Rx_2+Rx_3\subsetneq\dots,$$
     a contradiction.

     ($\Leftarrow$) Suppose that we have a chain
     $$M_1\subsetneq M_2\subsetneq M_3\subsetneq\dots,$$
     then $N=\bigcup_{i=1}^\infty M_i$
     is finitely-generated.
     Pick generators $x_1,\dots,x_r\in N$.
     Clearly $x_i$ has to be in some $M_j$'s,
     so one can pick $n$ such that $x_1,\dots,x_r\in M_n$,
     so $N=M_n$, and the chain terminates.
\end{proof}
Now, we start giving how to generate Noetherian modules.
\begin{proposition}[Exact Sequence Property]
    \label{prop:chain_condition_exact}
    Consider an exact sequence of $R$-modules
    $$0\longto M_1\stackrel\phi\longto M_2\stackrel\pi\longto M_3\longto 0.$$
    Then,
    \begin{enumerate}[label=(\alph*)]
    \item $M_2$ Noetherian if and only if $M_1$ and $M_3$
	are Noetherian.
    \item $M_2$ Artinian if and only if $M_1$ and $M_3$ 
	are Artinian.
    \end{enumerate}
\end{proposition}
\begin{proof}
    \begin{enumerate}[label=(\alph*)]
    \item ($\Rightarrow$)
    An infinite chain of submodules of $M_1$
    $$P_1 \subsetneq P_2 \subsetneq \dots \subsetneq M_1$$
    turns to an infinite chain of submodules in $M_2$:
    $$\phi(P_1) \subsetneq \phi(P_2) \subsetneq \dots \subsetneq M_2$$
    (they are distinct since $\phi$ is injective).
    Thus, if $M_1$ is not Noetherian, then $M_2$
    is not Noetherian.
    
    Similarly, an infinite chain of submodules of $M_3$
    $$Q_1 \subsetneq Q_2 \subsetneq \dots \subsetneq M_3$$
    turns to an infinite chain of submodules in $M_2$:
    $$\pi^{-1}(Q_1) \subsetneq \pi^{-1}(Q_2) \subsetneq \dots \subsetneq M_2$$
    (they are distinct since $\pi$ is injective).
    Thus, if $M_3$ is not Noetherian, then $M_2$
    is not Noetherian. Hence, we are done.
    
    ($\Leftarrow$)
    Given $N_1\subset N_2\subset N_3\subset\dots$
    is a chain in $M_2$.
    Then, if $\pi : M_2\to M_3$ again,
    then we have two chains
    $$
    \begin{array}{cccccl}
        \phi^{-1}(N_1)& \subset & \phi^{-1}(N_2)
        & \subset &  \phi^{-1}(N_3)
        & \subset\dots \\[4pt]
        \pi(N_1) & \subset  & \pi(N_2) & \subset &  \pi(N_3) & \subset\dots.
    \end{array}
    $$
    Since $M_1$ and $M_3$ are Noetherian,
    both chains stabilizes, so for some 
    sufficiently large $k$, we have
    $\phi^{-1}(N_k) = \phi^{-1}(N_{k+1})$
    and $\pi(N_k)=\pi(N_{k+1})$.
    We claim that this implies $N_k=N_{k+1}$.
    
    To see this, suppose $x\in N_{k+1}$, then
    $\pi(x)\in \pi(N_{k+1}) = \pi(N_k)$, 
    so for some $y\in N_k$, we have
    $\pi(x-y)=0$, which means $x-y\in \Im\phi$.
    Thus, there exists $z\in M_1$ such that $x=y+\phi(z)$.
    Now, $\phi(z)\in N_{k+1}$, so $z\in \phi^{-1}(N_{k+1})$,
    which implies $z\in \phi^{-1}(N_k)$, so $\phi(z)\in N_k$.
    Hence, $x = y+\phi(z) \in N_k$.
    \item is similar to (a). \qedhere
\end{enumerate}
\end{proof}
\begin{corollary}
    If $M_1,\dots,M_n$ are all Noetherian (resp. Artinian),
    then $M_1\oplus \dots\oplus M_n$ is Noetherian
    (resp. Artinian).
\end{corollary}
If $R$ is Noetherian ring,
checking if a module is Noetherian
is even more convenient.
\begin{proposition}[Noetherian iff finitely-generated]
    If $R$ is Noetherian ring,
    then $M$ is Noetherian if and only if 
    $M$ is finitely-generated as an $R$-module.
\end{proposition}
\begin{proof}
    $(\Rightarrow)$ is clear from \Cref{prop:noetherian_module_fin_gen}.
    For $(\Leftarrow)$, consider an exact sequence
    $R^{\oplus n}\to M\to 0$.
    We know that $R^{\oplus n}$ is Noetherian,
    so $M$ is Noetherian by the exact
    sequence property \Cref{prop:chain_condition_exact} (a).
\end{proof}
\begin{exercise}
Let $R$ be a Noetherian ring.
\begin{alphenum}
\item Prove that any submodule of finitely-generated $R$-module 
is also finitely-generated.
\item Prove that any finitely-generated module is 
finitely-presented.
\end{alphenum}
\end{exercise}
\begin{exercise}
    If $R$ is an Artinian ring, prove that 
    $M$ is Artinian module if and only if $M$ is
    finitely-generated as an $R$-module.
\end{exercise}
%
%
%
%
%
\section{Noetherian Rings}
\label{sec:noetherian}
\subsection{Properties of Noetherian Rings}
We now restrict our attention to Noetherian rings.
First, from \Cref{prop:noetherian_module_fin_gen}, we have
\begin{corollary}
$R$ is Noetherian
if and only if every ideal of $I$ is finitely-generated.
\end{corollary}
Now, we begin stating properties.
\begin{proposition}[Quotient, localization, product
of Noetherian ring is Noetherian]
    \label{prop:noetherian_ring_properties}
    Let $R$ be a Noetherian ring.
    \begin{enumerate}[label=(\alph*)]
    \item $R/I$ is Noetherian for all ideal $I$;
    \item $S^{-1}R$ is Noetherian for any multiplicative set $S$;
    \item $R\times S$ is Noetherian for any Noetherian rings
    $R$ and $S$.
    \end{enumerate}
\end{proposition}
\begin{proof}
    \begin{enumerate}[label=(\alph*)]
    \item
    Let $\pi : R\to R/I$ be the quotient map.
    By correspondence theorem \Cref{thm:correspondence},
    all ideals of $R/I$ must be in form $\pi(J)$ 
    for some ideal $J\subseteq R$.
    If $J = (x_1,\dots,x_n)$,
    then $\pi(J) = (\pi(x_1),\dots,\pi(x_n))$ is finitely-generated.
    \item
    By \Cref{prop:ideal_is_extended},
    any ideal $J\subset S^{-1}R$ is an extension
    of some ideal $I\subset R$.
    If $I=(x_1,\dots,x_n)$, then
    $J = \left(\frac{x_1}1, \dots,\frac{x_n}1\right)$.
    \item 
    Let $\pi_1 : R\times S\to R$ and $\pi_2 : R\times S\to S$
    be the projection maps.
    Given an ideal $J\subset R\times S$,
    let $J_R = \pi_1(J)$ and $J_S = \pi_2(J)$.

    We claim that $J = J_R\times J_S$.
    To do so,
    given $(a,b)\in J_R\times J_S$.
    We have $(a,c)\in J$ for some $c$
    and $(d,b)\in J$ for some $d$, then
    $$(a,b) = (a,c)\cdot (1,0) + (d,b)\cdot (0,1)\in J.$$
    Thus, $J$ is finitely-generated.
    \qedhere
    \end{enumerate}
\end{proof}
\begin{theorem}[Hilbert's Basis Theorem]
    \index{Hilbert's basis theorem}
    \label{thm:hilbert_basis}
    If $R$ is Noetherian, then so is $R[x]$
    (and hence $R[x_1,\dots,x_n]$).
\end{theorem}
This allows us to deduce, for example, that
$\mathbb C[x_1,\dots,x_n]$ and
$\mathbb Z[x_1,\dots,x_n]$ are Noetherian.
This is a nontrivial result,
since ideals $\mathbb C[x_1,\dots,x_n]$
can be quite arbitrary, but it is still finitely-generated.
It is also interesting to find an algorithm
finding these finitely-many generators.
However, typical proofs (e.g., the one below)
is not algorithmic in nature.
The set of generators you get is exponential.
\begin{proof}
    Let $I\subset R[x]$ be an ideal.
    We want to show that $I$ is finitely-generated.
    We want to use the fact that $R$ is Noetherian.
    The key idea is to look at the leading coefficients.
    Define an ideal
    $$J_n = \{a\in R : \text{there exists } f\in I
	\text{ with degree } n\text{ and leading 
    coefficient } a\} \cup \{0\}.$$
    Clearly, $J_n$ is an ideal.
    Moreover, $J_n\subset J_{n+1}$ because if $f\in I$
    of degree $n$, then $xf\in I$ of degree $n+1$.
    Therefore, since $R$ is Noetherian,
    we have $J_n=J_{n+1}=\dots$ for all $n$.
    Even further, for each $m\leq n$,
    we have $J_m$ is finitely generated, say by
    $(a_{m,i})$. We let $a_{m,i}$ be
    the leading coefficient of $f_{m,i}$.

    We now claim that $I = (f_{m,i} : m\leq n)$
    (and henceforth finitely-generated).
    To prove this, we prove by induction on $r=\deg f$
    that $f$ is generated by $f_{m,i}$'s.
    The base case $r=0$ is clear.
    We have two cases.
    \begin{itemize}
    \item If $r\geq n$, let $f(x) = cx^r + \cdots$.
	Then, $c\in J_r = J_n$, so $c$ is generated
	by $(a_{n,i})$. This means there exists
	$b_i\in R$ such that $c = \sum b_i a_{n,i}$.
	Then, we have
	$$\sum b_i f_{n,i}(x)\cdot x^{r-n} = cx^r
	+ \text{lower-order terms}.$$
	The left-hand-side is generated by $f_{n,i}$,
	so one can use the induction hypothesis.
    \item If $r\neq n$, the same thing happen:
	just change $f_{n,i}$ to $f_{r,i}$. \qedhere
    \end{itemize}
\end{proof}
To summarize, any quotient and localization
of Noetherian ring is Noetherian.
Adjoining finitely many elements into
Noetherian ring is also Noetherian.
This pretty much covers everything you will care about
in Algebraic Geometry.

We end with some warning.
\begin{warn}    
    \begin{itemize}
    \item If $R$ is a Noetherian ring,
    then subring $S\subseteq R$ is not necessarily Noetherian.
    Example: $S=k[x_1,x_2,\dots]$ and $R=\Frac S$.
    \item If $R$ and $S$ are $T$-algebra, 
    and $R$, $S$, $T$ are all Noetherian,
    then $R\otimes_T S$ is not necessarily Noetherian.
    The counterexample is $\QQ(x_1,x_2,\dots) \otimes_\QQ 
    \QQ(x_1,x_2,\dots)$ is not Noetherian;
    see \cite[Theorem 11]{Vamos}
    \end{itemize}
    Non-Noetherian rings show up in number theory
    (specifically, class field theory):
    for example, the adele ring $\bbA_\QQ$ is not Noetherian.
\end{warn}
%
%
%
%
\subsection{Zariski Topology of Noetherian Rings}
\label{subsec:irreducible_comp}
We have discussed a lot on properties on Noetherian rings.
Now, we try to interpret this in terms of spectrum.
\begin{definition}[Noetherian Space]
    A topological space $X$ is \vocab[Noetherian topological space]{Noetherian}
    if the closed subsets satisfy the Descending Chain
    Condition (DCC).
    In particular, there is no infinite chain
    of closed sets
    $$X\supset Z_1\supsetneq Z_2\supsetneq Z_3\supsetneq\dots.$$
\end{definition}
Most topological spaces are not Noetherian though.
For example, the real line $\mathbb R$ is obviously
not Noetherian because one can take infinitely nested intervals.

There is no surprising that this definition
is rigged so that it is equivalent to Noetherian rings.
\begin{proposition}
    If $R$ is Noetherian, then $\Spec R$ is a Noetherian space.
\end{proposition}
\begin{proof}
    There exists a correspondence between
    \begin{align*}
        \{\text{closed sets of }\Spec R\}
        &\ \leftrightarrow\ \{\text{radical ideals of }R\} \\
        Z &\ \mapsto\ I(Z) = \bigcap_{\pp\in Z} \pp \\
        V(J) &\ \mapsfrom\ J
    \end{align*}
    Observe that $I(V(J)) = \sqrt{J}$ by definition of radical,
    so these correspondence are inverses.
    
    Thus, given a chain of closed sets
    $Z_1\supsetneq Z_2\supsetneq Z_3\supsetneq\dots$,
    we get an ascending chain
    $$I(Z_1)\subsetneq I(Z_2)\subsetneq \dots,$$
    which violates the Noetherian condition.
\end{proof}
\subsection{Irreducible Components}
We now give a way to decompose $\Spec R$.
\begin{definition}[Reducible topological spaces]
    A topological space $X$ is \vocab[reducible topological space]{reducible} if $X=Y_1\cup Y_2$
    (not necessarily disjoint)
    such that $Y_i$ is closed and $Y_i\neq X$.

    If $X$ is not reducible and $X\neq\emptyset$, we say that $X$
    is \vocab[irreducible topological space]{irreducible}.
\end{definition}
Let us make some examples for the spectrum.
\begin{example}
    \begin{itemize}
    \item $\Spec\mathbb Z$ is irreducible
    because closed set is finite,
    so no two closed sets union to $X$.
    \item More generally, if $R$ is a domain,
    then $\Spec R$ is irreducible, because if 
    $\Spec R = Y_1\cup Y_2$, then $(0)$ must be in at 
    least one of $Y_i$'s, but closure of $(0)$
    is the entire $\Spec R$.
    \item On the other hand, $\Spec\mathbb C[x,y]/(xy)$
    is reducible because 
    $\mathbb C[x,y]/(xy)$ is the union of $x$-axis
    and $y$-axis. In particular, 
    $$\Spec \CC[x,y]/(xy) = V(x)\cup V(y).$$
    \end{itemize}
\end{example}
If $X$ is a topological space, 
then a closed subset $W\subseteq X$ is irreducible 
if and only if the induced topology on $W$
makes $W$ irreducible.
By plugging in the definition of induced topology,
we find that $W$ is irreducible if and only if
``for all closed subsets $Y,Z$ such that $W\subseteq Y\cup Z$,
we have either $W\subseteq Y$ or $W\subseteq Z$''.
An irreducible closed set is \vocab[maximal irreducible closed set]{maximal}
if it is not contained in another irreducible closed set.

Every Noetherian topological space can be decomposed
to a finite union of irreducible closed sets.
The Noetherian condition will guarantee finiteness.
\begin{proposition}[Irreducible decomposition]
    \label{prop:irreducible_decomp}
    Let $X$ be a Noetherian topological space.
    Then,
    $$X = \bigcup_{i=1}^n Z_i,$$
    where $Z_i$ is a maximal irreducible closed set.

    Moreover, this decomposition is unique. 
    Consequently, there are finitely many maximal
    irreducible closed sets.
\end{proposition}
\begin{proof} The proof is in two parts.

    \textbf{Existence.}
    If $X$ is irreducible, we are done.
    Otherwise, $X=Y_1\cup Y_2$ where $Y_i\subsetneq X$ closed.
    If both $Y_1$ and $Y_2$ are irreducible,
    we are done.
    Otherwise, take $Y_1=Y_3\cup Y_4$, where $Y_3,Y_4\subsetneq Y_1$.
    This process continues.
    It must terminate in finitely many steps 
    because $X$ is noetherian.
    Thus, we have $X= Z_1\cup \dots\cup Z_n$ 
    for some irreducible closed sets $Z_1,\dots,Z_n$.

    We now check that $Z_i$ is maximal. 
    Assume that $Z_i\subseteq Y$ 
    for some irreducible closed set $Y$.
    Then, note that $Y\subseteq Z_1\cup\dots\cup Z_n$,
    so $Y\subseteq Z_j$ for some $j$.
    Thus, $Z_i\subseteq Z_j$ for some $j$, and we can 
    throw $Z_i$ out of the decomposition.

    \textbf{Uniqueness.}
    If $X=Y_1\cup\dots\cup Y_m$ is another decomposition,
    then $Y_1\subseteq Z_1\cup\dots\cup Z_n$,
    so $Y_1\subseteq Z_i$ for some $i$, so by maximality,
    $Y_1=Z_i$. Thus, every set in $Y_i$ is equal to set in $Z_j$
    for some $j$.
    Conversely, every set in $Z_j$ is equal to $Y_i$ for some $i$.
    Therefore, $\{Y_i\}$ and $\{Z_j\}$ are the same set.
\end{proof}
We now interpret this result algebraically.
First, we can characterize when $\Spec R$ is irreducible.
\begin{proposition}
    The following are equivalent.
    \begin{alphenum}
    \item $\Spec R$ is irreducible.
    \item $\Nil(R)$ is a prime ideal.
    \end{alphenum}
\end{proposition}
\begin{proof}
    This is a straightforward definition-chasing.
    For completeness, we provide a complete proof below.
    
    \textbf{(a) implies (b).}
    If $a,b\notin \Nil(R)$ such that $ab\in\Nil(R)$,
    so $(ab)^n=0$ for some $n$.
    We claim that $\Spec R = V(a)\cup V(b)$.
    To see this, given any prime ideal $\pp$,
    $(ab)^n=0$, so $a\in\pp$ or $b\in\pp$,
    so $\pp\in V(a)$ or $\pp\in V(b)$.
    Finally, note that $V(a)\neq \Spec R$
    because $a\notin\Nil(R)$, so by \Cref{prop:nilradical},
    there exists a prime ideal not in $V(a)$.
    Similarly, $V(b)\neq\Spec R$. Thus, $\Spec R$
    is not irreducible.

    \textbf{(b) implies (a).}
    We claim that the closure of $\Nil(R)$ is $\Spec R$.
    To see why, consider a closed set $V(I)$ containing $\Nil(R)$,
    so $I\subseteq \Nil(R)$, which in turn contains $\pp$
    for all prime ideal $\pp$. Thus, every prime 
    $\pp\in\Spec R$ is in $V(I)$.

    Hence, if any irreducible closed set contains $\Nil(R)$,
    it must be entire $\Spec R$.
    This shows that $\Spec R$ is irreducible.
\end{proof}
\begin{corollary}
    \label{cor:prime_ideal_irreducible_closed_sets}
    \begin{alphenum}
    \item There exists a bijective correspondance
    \begin{align*}
        \{\text{prime ideals of }R\} 
        &\ \leftrightarrow\ \{\text{irreducible closed sets of }R\} \\
        \pp &\ \mapsto\ V(\pp).
    \end{align*}
    \item The correspondence above takes minimal prime ideals
    to maximal irreducible closed sets of $R$.
    \end{alphenum}
\end{corollary}
\begin{proof}
    \begin{alphenum}
    \item If $I$ is a radical ideal, the following are equivalent:
    \begin{itemize}
        \item $V(I) = \Spec R/I$ is irreducible.
        \item $R/I$ is a domain (proof: $R/I$ is reduced).
        \item $I$ is a prime ideal.
    \end{itemize}
    \item Follows from that $V$ is inclusion reversing,
    so it takes minimal prime ideals to
    maximal irreducible closed sets. \qedhere
    \end{alphenum}
\end{proof}
Thus, we have deduce the following corollaries.
\begin{corollary}[Noetherian rings have finitely many minimal primes]
    \label{cor:finitely_minimal_primes}
    If $R$ is Noetherian ring, then there are finitely
    many minimal primes.
\end{corollary}
\begin{proof}
    Combine \Cref{prop:irreducible_decomp} (on $X=\Spec R$)
    and \Cref{cor:prime_ideal_irreducible_closed_sets}.
\end{proof}
\begin{corollary}
    If $R$ is Noetherian ring, then
    for every ideal $I\subset R$, there are
    finitely prime ideals minimal among those containing $I$.
\end{corollary}
\begin{proof}
    Apply the previous corollary to $R/I$.
\end{proof}
Let us now interpret these results.
Let $R$ be a Noetherian ring and $I$ be a radical ideal.
Then, $V(I) \simeq \Spec(R/I)$ can be written as
union of $V(\pp)$, where $\pp$ is a prime ideal
minimal among those containing $I$. This means
$$V(I) = \bigcup_{i=1}^n V(\mathfrak p_i)
= V\left(\bigcap_{i=1}^n\mathfrak p_i\right).$$
Therefore, we have
$$I = \sqrt{\bigcap_{i=1}^n\mathfrak p_i}
= \bigcap_{i=1}^n\mathfrak p_i$$
where $\mathfrak p_i$ prime ideals minimal among
those containing $I$.
Thus, any radical ideal can be written as 
intersection of prime ideals minimal among those containing $I$.
Moreover, this decomposition is unique!

In the next Chapter, \Cref{sec:primary_decomp},
we will generalize this to non-radical ideals.
%
%
%
%
%
\section{Artinian Rings}
\label{sec:artinian}
Recall that Artinian rings satisfy
the Descending Chain Condition: there is no chain
$$I_1\supsetneq I_2\supsetneq I_3\supsetneq\dots.$$
We will now give the main structure theorem
for Artinian rings.

First, we study Artinian local rings.
Recall that a ring has dimension $0$ if and only if
every prime ideal is maximal.
\begin{theorem}[Artinian local rings]
\label{thm:artinian_local_ring}
Let $R$ be a local ring with maximal ideal $\mm$. 
The following are equivalent.
\begin{alphenum}
    \item $R$ is Artinian.
    \item $R$ is Noetherian and has dimension $0$
    (i.e., has $\mm$ as the only prime ideal).
    \item $R$ is Noetherian and $\mm = \Nil(R)$.
    \item $R$ is Noetherian and 
    $\mm^N=0$ for some integer $N\geq 0$.
\end{alphenum}
\end{theorem}
This theorem tells us that a local ring $R$ is Artinian 
if and only if $\Spec R$ has one point.
\begin{proof}
    We first do the easy part of showing that (b), (c), and (d) are all equivalent.
    \begin{itemize}
    \item \textbf{(b) $\Leftrightarrow$ (c).}
    Follows immediately from the fact that 
    $\Nil(R)$ is intersection of prime ideals (\Cref{prop:nilradical}).
    \item \textbf{(c) $\Rightarrow$ (d).}
    Let $\mm$ be generated by $x_1,\dots,x_n$.
    Since $\mm=\Nil(R)$, all elements $x_1,\dots,x_n$
    are nilpotent, so there exists $e$ such that
    $x_1^e = x_2^e = \dots = x_n^e = 0$.
    Now, we claim that $\mm^{ne}=0$. To see this, note that
    each term in $\mm^{ne}$ is a multiple of monomials
    $x_1^{f_1} \cdots x_n^{f_n}$ where $f_1+\dots+f_n\geq ne$.
    This monomial is zero because $f_i\geq e$ for some $i$.
    \item \textbf{(d) $\Rightarrow$ (c).}
    From $\mm^N=0$, taking radical both sides gives $\mm=\Nil(R)$.
    \end{itemize}

    Finally, we move on to the difficult part of showing that
    (a) and (d) are equivalent.
    This is done in two steps.
    \begin{enumerate}
    \item \textbf{If $R$ is Artinian,
    then $\mm^N=0$ for some $N$.}
    Consider a descending chain $\mm\supset \mm^2 \supset\dots$.
    Since $R$ is Artinian,
    we have $\mm^N = \mm^{N+1}$ for some $N$.

    It's so tempting to apply Nakayama's Lemma (\Cref{cor:nakayama_local_ring}),
    and conclude that $\mm^N=0$.
    However, this \textbf{doesn't work},
    as Nakayama requires that $\mm^N$ is finitely-generated.
    
    Let us find a way to salvage this.
    Assume $\mm^N\neq 0$.
    Let 
    $$\Sigma = \{\text{ideals } J\text{ such that }
    \mm^N J\neq (0)\}.$$
    We have $\Sigma\neq\emptyset$ since $(1)\in\Sigma$.
    We then find $J\in\Sigma$ that is minimal
    (this is a finite process because $R$ is Artinian).
    Then, there exists $x\in J$ such that $\mm^N x\neq 0$.
    However, $\mm^N \cdot (x)\neq 0$, so since $J$ is minimal,
    we have that $J=(x)$!
    
    Furthermore, we note that
    $$\mm^N(\mm J) = \mm^{N+1}J = \mm^N J \neq 0 
    \implies \mm J\in \Sigma,$$
    so by minimality of $J$ again, we get $\mathfrak mJ = J$.
    However, $J$ is finitely generated because it
    is principal! Thus, we may use
    Nakayama's lemma (\Cref{cor:nakayama_local_ring})
    to get that $J=(0)$,
    and this is a contradiction.
    \item \textbf{If $\mm^N=0$ for some $N$,
    then $R$ is Artinian if and only if 
    $R$ is Noetherian.}

    We consider the filtration
    $$R\supset \mathfrak m\supset \mathfrak m^2
    \supset \dots \mathfrak m^{N-1}\supset \mathfrak m^N=0.$$
    Then, the following are equivalent.
    \begin{itemize}
    \item $R$ is Noetherian as $R$-module.
    \item $\mm^{i-1}/\mm^i$ is Noetherian as $R$-modules for all $i=1,2,\dots,N$. (Use \Cref{prop:chain_condition_exact}.)
    \item $\mm^{i-1}/\mm^i$ is Noetherian as $R/\mm$-modules for all $i=1,2,\dots,N$.
    \item $\dim_{R/\mm}(\mm^{i-1}/\mm^i) < \infty$
    for all $i=1,2,\dots,N$.
    \item $\mm^{i-1}/\mm^i$ is Artinian as $R/\mm$-modules
    for all $i=1,2,\dots,N$
    (as an $R/\mm$-module).
    \item $\mm^{i-1}/\mm^i$ is Artinian as $R$-modules for all $i=1,2,\dots,N$
    (as an $R$-module).
    \item $R$ is Artinian (as $R$-module). 
    (Use \Cref{prop:chain_condition_exact} again.) \qedhere
    \end{itemize}
    \end{enumerate}
\end{proof}

Now, we can prove the main result for general Artinian ring.
\begin{theorem}[General Artinian rings]
    \label{thm:artinian}
   Let $R$ be a ring. The following are equivalent.
   \begin{alphenum}
    \item $R$ is Artinian.
    \item $R$ is Noetherian of dimension $0$.
    \item $\Spec R$ is finite and has discrete topology.
    \item $R$ is a finite product of Artinian local rings.
   \end{alphenum}
\end{theorem}
The moral of this theorem is that 
\begin{center}
\fbox{Artinian rings look like finitely many points.}
\end{center}
\begin{proof}
We first show that (b), (c), and (d) are equivalent.
\begin{itemize}
\item \textbf{(b) $\Rightarrow$ (c).} \Cref{cor:finitely_minimal_primes} gives that 
there are finitely minimal primes,
and from $R$ has dimension $0$, 
every prime ideal is a minimal prime,
so there are finitely many prime ideals.
\item \textbf{(c) $\Rightarrow$ (d).} By
\Cref{prob:product_ring}, if $\Spec R$
is a disjoint union of points, it is a product ring
$R_1\times\dots\times R_n$,
where $\Spec R_i$ is a single point. 
Note that $\Spec R_i$ is a single point, so by \Cref{thm:artinian_local_ring}
((b) implies (a)), we get that $R_i$ is an Artinian 
local ring.
\item \textbf{(d) $\Rightarrow$ (b).}
The dimension $0$ part follows from 
that the topology is discrete,
while the Noetherian part follows from 
\Cref{thm:artinian_local_ring} that every Artinian local ring 
is Noetherian and \Cref{prop:noetherian_ring_properties}.
\end{itemize}

Next, we show that (d) implies (a).
To do this, it suffices to show that product
of two Artinian rings is Artinian
(hence by induction, finite product of Artinian ring 
is Artinian).
Let $A$ and $B$ be Artinian. Then, we have an exact
sequence of $(A\times B)$-module
$$0\longto A\longto A\times B\longto B\longto 0.$$
Note that $A$ is Artinian as
$A$-module, so $A$ is Artinian as $(A\times B)$-module.
Simiarly, $B$ is Artinian as $(A\times B)$-module,
so by exact sequence property, \Cref{prop:chain_condition_exact}
$A\times B$ is Artinian.

Finally, we show the hardest implication that (a) implies (c).
This is in two steps.
\begin{enumerate}
\item \textbf{Every prime ideal of $R$ is maximal
(i.e., $R$ has dimension $0$).}
Let $\mathfrak p\subset R$ be prime but not maximal.
Then, $A=R/\mathfrak p$ is a domain but not a field.
Thus, there exist $0\neq f\in A$
such that $f$ is not a unit,
producing an infinite descending chain
$$A\supsetneq (f)\supsetneq (f^2)\supsetneq (f^3)
\supsetneq \dots,$$
which contradicts that fact that $A$ is Artinian.
\item \textbf{$R$ has finitely many maximal ideals.}
Suppose $\mm_1,\mm_2,\dots$ are distinct maximal ideals of $R$.
Then, we product a chain 
$$\mm_1\supsetneq \mm_1\mm_2 \supsetneq \mm_1\mm_2\mm_3
\supsetneq\dots$$
Note that equality cannot occur:
if $\mm_1\cdots\mm_n = \mm_1\cdots\mm_{n+1}$,
then $\mm_1\cdots\mm_n\subseteq \mm_{n+1}$,
so the primeness of $\mm_{n+1}$ implies $\mm_i\subseteq\mm_{n+1}$
for some $i$, which is a contradiction.
Thus, we have an infinite descending chain.
\qedhere
\end{enumerate}
\end{proof}
\begin{example}
    Let $R=\mathbb C[x_1,\dots,x_n]/I$,
    and suppose that the only point in $\Spec R$
    is $\vec 0$. Thus, $(x_1,\dots,x_n)$ is
    going to be the unique maximal ideal,
    and $\mathfrak m^k=0$ in $R$ for some $k$.
    Therefore, if $I$ is an ideal satisfying
    $$(x_1,\dots,x_n)^k \subset I\subset (x_1,\dots,x_n),$$
    then $\mathbb C[x_1,\dots,x_n]/I$ is Artinian.
    For example, the ring
    $$\mathbb C[x]/(x^3, x^2y, xy, xy^5)$$
    is Artinian.
\end{example}

\section{Problems}
\begin{problem}
\begin{alphenum}
\item Let $A$ be a Noetherian ring and $\phi : A\to A$ be a surjective ring morphism. Prove that $\phi$ is an isomorphism.

(Hint: consider $\Ker\phi^n$.)
\item Let $A$ be a ring, and let $M$ be a Noetherian module.
Prove that every surjective module morphism $\phi : M\to M$
is an isomorphism.
(Notice that this is a special case of \Cref{prob:surjective_map}.
You are asked to come up with a different proof.)
\item Prove that (a) and (b) holds when ``Noetherian'' 
is replaced with ``Artinian'' and ``surjective''
is replaced with ``injective''.
\end{alphenum}
\end{problem}
\begin{problem}[Cohen's Theorem]
Prove that a ring $A$ is Noetherian if and only if 
every \textbf{prime} ideal of $A$ is finitely-generated.
\end{problem}
\begin{problem}
\label{prob:algebraic_integer}
Prove that the ring of algebraic integer (i.e.,
the integral closure of $\ZZ$ in $\overline\QQ$)
is not Noetherian.
\end{problem}
\begin{problem}
\label{prob:formal_series}
Let $A$ be a Noetherian ring. Prove that the formal series ring
$$A[[x]] := \{a_0+a_1x+a_2x^2+\dots \,:\, a_0,a_1,\dots,\in A\}$$
is Noetherian.
\end{problem}
\begin{problem}
Let $A$ be a ring such that $A_\pp$ is Noetherian 
for all prime ideal $\pp$. 
Must $A$ be Noetherian?
\end{problem}
\begin{problem}
Let $A$ be a ring such that $A$ is Noetherian, 
$\operatorname{Spec} A$ is connected, 
and $A_{\mathfrak p}$ is a domain for all prime ideal 
$\mathfrak p$. Prove that $A$ is a domain. 
\end{problem}
\begin{problem}
Let $k$ be a field and $A$ be a finitely-generated $k$-algebra.
Prove that $A$ is Artinian 
if and only if $\dim_k A$ is finite. 
(There is a very short solution.)
\end{problem}